\scanpage{197}
\mspage{II.45}
% archive sheet number: 196
% lines 1--6
Il y a lieu de travailler dans la catégorie $\underset{\lambda}{\varinjlim}$ (faisceaux \unsure{cts} tordus sur $X_\lambda$), qui précisément s'explicite en termes de représentations \uline{linéaires} \inserted{$\unsure{\ell|l}$-adiques} du pro groupe $\Pi$, ou en termes de $\Pi^{\unsure{\mathrm{alg}|\mathrm{an}}}$, ou de rep. alg. de $\Pi^{\unsure{\mathrm{alg}|\mathrm{an}}}$.

% lines 7--13
\uline{Exemple 1} Soit $f:Y\unsure{-|\to}X$ un morphisme propre et lisse, on prend les $R^i f_*(\deleted{\unsure{\mathbb{Z}}}\mathbb{Q}_{\ell Y})$, \unsure{on|ou} \deleted{\unsure{en déduit|on déduit}} les $R^i f_*(\deleted{\unsure{\mathbb{Z}}}\mathbb{Q}_{\ell Y})$ si on veut une représentation de $\unsure{\widetilde G|\widetilde C},\unsure{\widetilde{\mathfrak g}|\widetilde g}$. De $=$, les hom entre $R^i f_*$ définis par des $X$-morphismes $Y\unsure{-|\to}Y'$, ou $[$\deleted{\unsure{si}} $X$ est intègre normal) par des \deleted{classes de} correspondances \inserted{algébriques} sur les fibres \unsure{génériques|géométriques}.

% lines 14--20
\uline{Exemple} $^2$ (\unsure{ess\textsuperscript{t} identique}). $T_\ell(A)$, où $A$ est un schéma abélien sur $X$, \deleted{$\unsure{=T_\ell(A)\otimes_{\mathbb{Z}_\ell}\mathbb{Q}_\ell}$} Noter que $\underline{\operatorname{Hom}}(T_\ell(A),\deleted{\unsure{\mathbb{Q}}}\mathbb{Z}_\ell)\simeq R^1 f_*(\mathbb{Z}_{\ell A})$, donc $T_\ell(A)$ et $R^1 f_*(\mathbb{Z}_{\ell A})$ se déterminent mutuellement. \deleted{\illegible{2}} Si on veut en déduire une représentation de $\widetilde{\mathfrak g}=\operatorname{Lie}(\pi_1^{\mathrm{arith}}(X;\unsure{\xi|\eta}))$, il faut \unsure{tensoriel|tensoriser} d'abord par $\mathbb{Q}_\ell$.
